%% Согласно ГОСТ Р 7.0.11-2011:
%% 5.3.3 В заключении диссертации излагают итоги выполненного исследования, рекомендации, перспективы дальнейшей разработки темы.
%% 9.2.3 В заключении автореферата диссертации излагают итоги данного исследования, рекомендации и перспективы дальнейшей разработки темы.
\begin{enumerate}
  \item На основе анализа \ldots
  \item Численные исследования показали, что \ldots
  \item Математическое моделирование показало \ldots
  \item Для выполнения поставленных задач был создан \ldots
\end{enumerate}


Основные результаты работы заключаются в следующем.
1. Произведен обзор методов переноса знаний при решении проблем обра
ботки естественного языка.
2. Произведено внутреннее и внешнее сравнение контекстно-независи
мых векторных представлений слов, обученных разными способами
на ограниченной обучающей выборке для трех малоресурсных языков.
Обученные векторные представления выложены в открытый доступ.
3. Разработана вопросно-ответная модель на основе BERT с применением
библиотеки transformers от HuggingFace.
4. Показано, что межъязыковой перенос позволяет использовать обще
доступные обучающие данные английского языка при дообучении
M-BERT для вопросно-ответной задачи, таким образом, отсутствует
необходимость в сборе полноценного обучающего датасета целевого
языка.
5. Экспериментально установлено, что применяя метод межъязыково
го переноса, M-BERT, дообученный с применением многоязычной
обучающей выборки, имеет сопоставимое качество с языко-специ
фичными BERT для вопросно-ответной задачи, тем самым отпадает
необходимость в вычислительных ресурсах для предобучения языко
специфичных BERT.
6. Разработана оригинальная модель GOLOMB, которая использует во
просно-ответную модель SQuAD для отслеживания состояния диалога.
7. Модели, обученные в рамках работы, доступны в открытом доступе.